% Несъгласие за границите на една заявка по HTTP/1.1.
%
% Източник: src/core/request_parse.cpp (parse_request_head) и
% include/coroute/core/request.hpp (fold, add_header, header).
%
% Рисува се като една лента от октети с два четеца отгоре и отдолу. Причината да
% не се рисуват две отделни заявки е, че октетите са едни и същи: спорът е само
% къде свършва първата. Затова лентата е обща, а разликата се носи от скобите
% над и под нея.
%
% Черно на бяло: преденият четец е с плътна скоба, този сървър е с прекъсната.
% Спорните октети са щриховани, а не оцветени. Никакъв цвят.
\begin{figure}[htbp]
\centering
\begin{tikzpicture}[
    % Кутия за един случай и кутия за поправката му. Поправката е със сив фон,
    % за да се чете като друг слой: какво беше срещу какво е сега.
    % Еднаква минимална височина за трите случая, за да застанат поправките под
    % тях на един ред. Иначе трите кутии свършват на различна височина.
    casebox/.style={box, align=left, text width=3.95cm, minimum height=2.45cm,
                    font=\fontsize{7.6}{9.2}\selectfont, inner sep=4pt},
    fixbox/.style={casebox, fill=gray!25, minimum height=1.75cm},
    seg/.style={draw, thick},
    hdrline/.style={anchor=north west, font=\fontsize{8}{9.6}\selectfont,
                    inner sep=1pt},
    brc/.style={decorate, decoration={brace, amplitude=5pt}, thick},
  ]

  % ---- лентата от октети --------------------------------------------------
  % Три части: заглавна част, спорни октети, следващи октети във входа.
  \draw[seg] (0,1.2) rectangle (6.4,3.4);
  \draw[seg, pattern=north east lines] (6.4,1.2) rectangle (9.6,3.4);
  \draw[seg, fill=gray!10] (9.6,1.2) rectangle (13.8,3.4);

  % Заглавната част е изписана ред по ред, за да има къде да седне полето за
  % дължина. То е нарисувано като гнездо с прекъсната рамка: трите случая долу
  % са трите начина то да бъде запълнено.
  \node[hdrline] at (0.15,3.32) {\code{POST /p HTTP/1.1}};
  \node[hdrline] at (0.15,2.96) {\code{Host: h}};
  \node[draw, densely dashed, thick, fill=gray!25, rounded corners=1pt,
        font=\fontsize{8}{9.6}\selectfont, inner sep=3pt, anchor=north west]
       at (0.12,2.56) {поле за дължина: вж. (1), (2), (3)};
  \node[hdrline] at (0.15,1.92) {\code{CRLF CRLF}, край на заглавната част};

  % Надписът върху щриховката е с бял фон, иначе не се чете.
  \node[font=\fontsize{8}{9.6}\selectfont, align=center, fill=white, inner sep=2pt]
       at (8.0,2.30) {8 октета\\спорни};

  \node[font=\fontsize{8}{9.6}\selectfont, align=center] at (11.7,2.30)
       {\code{GET /admin ...}\\следващи октети\\във входа};

  % ---- четец 1: преденият сървър ------------------------------------------
  \draw[brc] (0,3.55) -- (9.6,3.55);
  \node[note, anchor=south, align=center] at (4.8,3.85)
       {преден сървър: заявката заема тези октети};

  % ---- четец 2: този сървър -----------------------------------------------
  % Скобите се чертаят отдясно наляво, за да изпъкнат надолу. Прекъснатата
  % линия ги отличава от плътната скоба на предния четец. Малкият процеп при
  % 6.4 е само за да не се слеят двата върха в една фигура.
  \draw[brc, densely dashed] (6.3,1.05) -- (0,1.05);
  \draw[brc, densely dashed] (13.8,1.05) -- (6.5,1.05);
  \node[note, anchor=north, align=center] at (3.2,0.65)
       {този сървър: заявката свършва тук};
  \node[note, anchor=north, align=center] at (10.1,0.65)
       {остатък: за този сървър тук започва\\следващата заявка};

  % ---- трите източника на неяснотата ---------------------------------------
  \node[casebox, anchor=north west] (c1) at (0,-0.95)
       {\textbf{(1)} Малки букви в името\\
        \code{content-length: 8}\\[2pt]
        преден сървър: тяло от 8 октета\\
        този сървър: търсене с точна\\
        главност не намира поле,\\
        тялото е 0 октета};
  \node[casebox, anchor=north west] (c2) at (4.7,-0.95)
       {\textbf{(2)} Две дължини, които спорят\\
        \code{Content-Length: 8}\\
        \code{Content-Length: 0}\\[2pt]
        преден сървър взима първата\\
        речник по име пази последната\\
        8 октета срещу 0 октета};
  \node[casebox, anchor=north west] (c3) at (9.4,-0.95)
       {\textbf{(3)} Подминато поле\\
        \code{Transfer-Encoding:}\\
        \code{chunked}\\[2pt]
        преден сървър чете тялото\\
        четец, който подминава полето,\\
        не чете нито един негов октет};

  \node[fixbox, anchor=north west] at (0,-3.52)
       {Поправено: имената се сгъват по ASCII и се търсят сгънати.
        RFC 9110, 5.1};
  \node[fixbox, anchor=north west] at (4.7,-3.52)
       {Поправено: две различни стойности се отказват с 400.
        RFC 9112, 6.3};
  \node[fixbox, anchor=north west] at (9.4,-3.52)
       {Поправено: само \code{chunked} се декодира. Друго кодиране: 501.
        Заедно с \code{Content-Length}: 400. RFC 9112, 6.1 и 6.3};

  % ---- твърдението, което фигурата поддържа --------------------------------
  \node[grp, note, text width=13.2cm, fill=gray!10, anchor=north west] at (0,-5.50)
       {И трите вече се отказват или се обработват. Дефектът не е в отделното
        поле, а в несъгласието: двата четеца броят различен брой октети за една
        и съща заявка. При този сървър остатъкът се изхвърля заедно с буфера за
        четене, така че само по себе си това не е контрабанда на заявки. То е
        предпоставката за нея.};

\end{tikzpicture}
\caption{Две реализации, които не се разбират къде свършва една заявка по
\code{HTTP/1.1}. Лентата е един и същ вход. Преденият сървър приписва спорните
октети на тялото, а този сървър ги смята за начало на следващата заявка. Трите
случая отдолу са трите начина полето за дължина да стане двусмислено, и трите
водят до една и съща картина. Поправките ги отказват или ги обработват.}
\label{fig:header_framing}
\end{figure}
